The study presents a Digital Twin framework for risk assessment in construction scheduling, combining traditional PERT/CPM techniques with Monte Carlo simulation and Bayesian Networks. The framework uses probabilistic analysis based on financial risk management metrics, such as Value-at-Risk (VaR) and Conditional Value-at-Risk (CVaR), to evaluate uncertainty in project duration.

The model was implemented in a Python-based client-server system with a Streamlit interface, allowing updates of project duration estimates and providing a more adaptive approach for construction planning under uncertainty.

The paper contributes to the literature by combining concepts from construction scheduling, probabilistic modeling, and Digital Twin systems. The Bayesian Network structure allows updating predictions when new evidence is available. From a practical perspective, the framework can support project managers in risk assessment and monitoring activities.

However, the study has relevant limitations. The framework was tested using a simplified toy problem, which may not represent the complexity of real construction projects. In addition, activity durations were modeled using triangular distributions, which may simplify real uncertainties and dependencies. Furthermore, the proposed framework operates more as a digital shadow than a fully integrated digital twin.

Future studies could apply the framework to real construction projects, integrate real-time sensing technologies, and explore more advanced inference and automation mechanisms.
